
The quadrupole tensor dependence in the magnetic field direction can also be obtained from the pure Kitaev model when we consider the second-order perturbation theory developed in the Appedix~\ref{App-pertubation}.
The perturbation at second order renormalizes the Kitaev exchanges in \eqref{eq:H0spin} to generate anisotropic ones $$K_\gamma-K \propto -(h_\gamma)^2$$
with proportionality constant going with the inverse of the vison gap $ \sim 1/\Delta_{2\mathrm v}$. 


%In a first approximation, let us discuss the dependence of the quadrupole tensor on the magnetic field direction by treating the field perturbatively in the pure Kitaev model.  For magnetic field directions not perpendicular to the lattice plane, the (often discarded) contribution from second-order perturbation theory generates an anisotropic renormalization of the Kitaev couplings. This effect is captured by the Hamiltonian in Eq.~(\ref{eq:pure_K}) with 


In Fig.~\ref{fig:quadrupole-pure}, we show the angular dependence of the quadrupole components  $Q_{33}$, $Q_{11}-Q_{22}$ and $Q_{12}$ calculated from the spin correlations for the model in Eq.~(\ref{eq:pure_K}). 
The component $Q_{33}$ does not change sign but varies slightly around an average value with an angular dependence qualitatively similar to  $|h_xh_yh_z|$. In particular, $Q_{33}$ is the maximum for a field along the $\hat{\mb c}$ direction, which may be interesting to maximize the intrinsic electric field produced at positions right above the vortex.    

\begin{figure}[t]
    \hspace*{-10mm}%
    \includegraphics{Figures/npj/fig6.pdf}
    \caption{Quadrupole components as a function of magnetic field direction, calculated using the exactly solvable  Hamiltonian with coupling generated by second and third order perturbation theory in the magnetic field. The coupling constants $K_\gamma$ and $\kappa$ were calculated using $\kappa = 0.338  \; \frac{h_xh_yh_z}{\Delta_{2\mathrm v}^2}$ and $K_\gamma=K-   \frac{(h_\gamma)^2}{\Delta_{2\mathrm v}}$ with  $|\mathbf{h}|=0.2 |K|$ and $\Delta_{2\rm v} = 0.263 |K|$.
    The scale is in units of $t_2^2t_2^\prime/U^3$ and we set the lattice spacing to unity.  Here we use $L=42$.}
    \label{fig:quadrupole-pure}
\end{figure}
On the other hand, the difference $Q_{11}-Q_{22}$ vanishes for $\mb h\parallel \hat{\mb c}$, but is maximum when the field points along the $\hat {\mb z}$ axis; this is the direction in which the anisotropy in the effective Kitaev couplings is maximized, with $K_z<K_x=K_y$.   Finally, $Q_{12}$ vanishes if we tilt the field along the high-symmetry $ac$ plane, but becomes nonzero for more general field directions.









